\documentclass[../main.tex]{subfiles}



\begin{document}

\section{Starting Point and SFT Data Preparation}
\label{sec:starting_point}

\subsection{The Mini-Assignment 1 Checkpoint}
\label{sec:ma1_checkpoint}

Mini-Assignment~1 delivered two variants of SmolLM2-360M: a full fine-tune and a LoRA adapter \cite{hu2022lora}. The LoRA variant was selected as the best in-domain fit and is the checkpoint carried forward here.

Before any further training, the adapter is merged into the base with \texttt{merge\_and\_unload}. Merging produces one consolidated checkpoint equivalent at inference time to (base + Mini-Assignment~1 LoRA), but with no PEFT bookkeeping. Every subsequent stage trains a fresh LoRA on top of a plain base, so there is no stacked-adapter arithmetic and no question of which adapter is active. The cost is a one-off 694~MB checkpoint on disk. From this point on the notebook treats the merged model as its pretraining starting point.

\subsection{SFT Corpus Generation}
\label{sec:sft_data}

A supervised fine-tuning dataset does not exist for this domain. It was generated by an ETL agent, \texttt{atlm\_teacher}, a Gemma~4~E2B model served locally via \texttt{agent\_server}. Given a raw Djinni posting the agent produces three short recruiter queries and one clean structured job description in Markdown (\texttt{\#\# Summary}, \texttt{\#\# Required Skills}, \texttt{\#\# Responsibilities}, \texttt{\#\# Requirements}), or a skip marker for unusable input. The output contract was validated on a diverse sample before the bulk run.

The bulk ETL converted 2{,}507 raw postings. Each record carries three independent recruiter queries pointing to the same job description; fanning these out yields 7{,}521 (query, posting) training pairs. Reuse of the same posting target under three different phrasings is a cheap augmentation that teaches the model phrasing-robustness without any extra annotation cost.

The records are split 90/10 into train and validation \emph{at the record level}, not at the query level, so no job description leaks across the split. With seed 42, this produces 2{,}257 training records (6{,}771 examples) and 250 validation records (750 examples). The split is also re-used in Section~\ref{sec:preference_dataset}: the 250 held-out validation records become the preference prompts, keeping the preference data strictly separate from SFT training without requiring additional annotation.

\subsection{Prompt Template}
\label{sec:prompt_template}

SmolLM2-360M has never been instruction-tuned; it has no notion of a request versus a response boundary. A consistent prompt template provides that signal during SFT and is reused unchanged at inference. The template is Alpaca-inspired:

\begin{verbatim}
You are a recruitment assistant. Given a brief recruiter
request, write a complete structured job posting in Markdown.

### Request
{query}

### Posting
{jd}
\end{verbatim}

At inference time the model receives everything up to and including \texttt{\#\#\# Posting\textbackslash n} and generates the posting from there. Three design decisions governed the choice over leaner key-value and ChatML alternatives:

\begin{enumerate}
  \item \textbf{Explicit task framing.} The one-sentence preamble costs roughly 25 tokens and gives the model an unambiguous anchor for the task.
  \item \textbf{No heading clash.} The structured postings already use \texttt{\#\#} headings; using \texttt{\#\#\#} for template separators lets the model distinguish template markers from response headings.
  \item \textbf{No special-token machinery.} Plain-text separators tokenise into existing vocabulary, so no new embeddings are added and no chat template needs registering on the tokenizer.
\end{enumerate}

\end{document}

\end{document}
